% =========================================================================
% DRAFT — Insertion-ready prose for Background Chapter 2 (2026-09-07)
% Two new sections proposed:
%   Section 2.9  State of the Art in Neuromechanical Simulation
%   Section 2.10 Testing Neuromechanical Hypotheses on Robots: The Ethical Case
% Place both AFTER current section 2.8 (Simulation Tooling) so the pipeline
% description precedes the field survey. Requires the 9 candidate references
% in Notes/candidate_references(.md/.bib) to be merged into thesis.bib first
% (keys: geyer_musclereflex_2010, song_deep_2021, degroote_perspective_2021,
%  webb_robots_2001, webb_robots_2002, russell_principles_1959,
%  tannenbaum_russell_2015, morton_cerebellar_2006, qin_vinsmono_2018).
% All citations verified 2026-09-07. Numbering will auto-adjust in Overleaf.
% =========================================================================

\section{State of the Art in Neuromechanical Simulation}\label{sec:bg_sota}

Neuromechanical simulation --- the coupling of a musculoskeletal plant to a model
of the neural controller that drives it --- has advanced along three partially
overlapping lines. The first line embeds reflex physiology directly into the
controller. The muscle-reflex walking model of Geyer and Herr
\citep{geyer_musclereflex_2010} demonstrated that positive force, length, and
velocity feedback around each muscle, organized with simple interlimb
coordination, is sufficient to produce self-stabilizing human-like walking
dynamics and realistic muscle activations without any central pattern generator.
That model became the reference point for a generation of reflex-based
neuromechanical controllers, including the cat hindlimb models of
Section~\ref{sec:bg_neural} \citep{markin_afferent_2010, shevtsova_modeling_2016}
and their extensions.

The second line replaces hand-designed circuitry with machine learning. Song and
colleagues \citep{song_deep_2021} trained deep reinforcement learning policies on
a physiologically realistic musculoskeletal model in OpenSim, reproducing human
kinematics and muscle activations across locomotion tasks, with the learned
policy re-usable for experiments that would be difficult in the laboratory. This
approach inherits the fidelity of the musculoskeletal plant but encodes control
in opaque function approximations: the trained policy can reproduce behavior, but
the individual ``neurons'' of the controller do not correspond to identifiable
physiological circuits, so the class of scientific questions it can answer ---
lesions, deletions, pathway stimulation --- is narrower than its behavioral
repertoire. Recent perspectives on predictive musculoskeletal simulation
\citep{degroote_perspective_2021} place both learned and optimal-control
approaches in the same frame: the plant models are mature, and the open problems
are the cost of solving them and the biological interpretability of the control
they find.

The third line, in which this dissertation sits, is biologically grounded
synthetic nervous systems: controllers whose neurons and synapses are explicit,
conductance-based, and mapped to identified physiological circuits, simulated at
speeds sufficient for real-time robotics \citep{nourse_sns_2023}. What
distinguishes the present program from all three lines is the combination it
attempts: an anatomically grounded controller (two-layer CPG with Ia, Ib, II,
and contact feedback), a biomechanically validated pneumatic plant
(Chapters~\ref{ch:methods}--\ref{ch:results}), and an eventual physical robot.
Simulation speed is the practical enabler: the MuJoCo backend evaluates the
musculoskeletal loop faster than real time on commodity hardware
\citep{todorov_mujoco_2012, caggiano_myosuite_2022}, and SNS-Toolbox sustains
real-time or faster network evaluation from workstations down to embedded boards
\citep{nourse_sns_2023}. A single overnight run can therefore sweep thousands of
parameter combinations or execute hundreds of simulated gait cycles per
condition --- throughput that no human-subject or animal protocol could match,
and that transforms lesion and stimulation studies from bespoke demonstrations
into statistically powered experiments.

\section{Testing Neuromechanical Hypotheses on Robots: The Ethical Case}\label{sec:bg_ethics}

The experiments this research program is designed to enable --- deleting a
muscle from a limb, blocking a sensory pathway, lesioning part of a rhythm
generator, applying tonic stimulation to spinal circuits --- cannot be performed
on healthy humans. Invasive intervention on the human neuromuscular system is
limited to clinical populations, is tightly constrained by ethics review, and
permits neither the within-subject control nor the reversibility that
neuromechanical hypotheses demand. The same experiments in animals carry real
moral cost, and their justification is governed by the Three Rs framework:
replacement of animal subjects with non-animal methods, reduction of the numbers
required, and refinement of the remaining procedures
\citep{russell_principles_1959, tannenbaum_russell_2015}. A robot that implements
a validated neuromechanical hypothesis is precisely such a replacement method:
every deletion, lesion, and stimulation experiment that runs on the synthetic
nervous system is an experiment that does not require a living subject, and the
robot population can be duplicated exactly, so that ``subjects'' are
statistically identical across conditions --- a form of reduction no breeding
colony can offer. Refinement follows as well: the failure modes of these
experiments damage hardware, not tissue.

For a robot to serve as a replacement subject, however, it must first earn
standing as a scientific model. Webb's influential analyses of biorobotic
methodology \citep{webb_robots_2001, webb_robots_2002} argue that a robot is a
good model of biological behavior when the mechanism it implements is tested at
the same level of organization as the target phenomenon --- matching dynamics,
control structure, and sensory ecology rather than surface appearance. The
validation chain of this dissertation is constructed to that standard: the
actuator force model is validated against bench data
(Section~\ref{sec:results_force}), the joint torque model against isometric
measurements and human torque envelopes (Section~\ref{sec:results_torque}), and
the neuromechanical controller will be verified against published tonic
stimulation and deletion phenomena \citep{selionov_tonic_2009,
rybak_modelling_2006} before any locomotion tuning is attempted. Only then can a
robot experiment be read as evidence about the hypothesized biological circuit
rather than about the robot.

The argument is ultimately bidirectional, completing the loop described in
Section~\ref{sec:project}. Experiments that are unethical, impractical, or
statistically unreachable in humans and animals become routine on the robot; the
resulting insights into neuromechanical control feed back into biology, into the
design of the next robot iteration, and directly into prosthetic and orthotic
devices that face the same problem of controlling muscles-like actuators for
human users. The ethical case and the scientific case are the same case: the
robot is the instrument that makes these experiments possible at all.
